% Chapter 9: Tables (practice author). PLAN.md section 1.4. Every body is generated by tools/generate_data.cpp from
% the library (data/table-*.tex, read with \tablebody) and checked there against the moment conditions.
% Tables: T4 to T7 tab:weights-central-{first,second,third,fourth}, T8 to T11 tab:weights-forward-{...},
% T12 tab:default-steps, T13 tab:ridders-constants.
% The macros \breakcommas, \weighttablestart and \weighttableend are in macros.tex (merged at stage 4).
\chapter{Tables}\label{ch:tables}
% Four weight tables may stack at the top of a page (main.tex restores the default of 3 after this chapter).
\setcounter{topnumber}{4}

The tables of this chapter are the single source of the weights and constants that Part~I quotes. The book's data
generator wrote the body of every table from the library itself, from the same internal tables that
\code{derivative} uses. The few constants that are literals of the library's code are transcribed instead, as
\cref{sec:tables-ridders} says. Before it writes a rule, the generator checks the rule against its moment conditions.
The library's tests compare the rules with the published tables of Fornberg\cite{fornberg1988} and Chapra
(Chapter 23).\cite{chapra2015}

\section{How to read the tables}\label{sec:tables-reading}

A rule of accuracy order $\aorder$ for the derivative of order $\dorder$ samples $f$ at the nodes $x + n_k h$. It
combines the samples with integer numerators $c_k$ over one common denominator $\den$:
\begin{equation}
  f^{(\dorder)}(x) \approx \frac{1}{\den\, h^{\dorder}} \sum_k c_k\, f(x + n_k h).
  \label{eq:tables-rule}
\end{equation}
The weights of \cref{ch:taylor,ch:stencils} are $w_k = c_k / \den$. They are fixed by the moment conditions
(\cref{sb:moments}). With the moments $\mom{j} = \sum_k c_k n_k^j$, the conditions read
\begin{equation}
  \mom{\dorder} = \dorder!\, \den, \qquad \mom{j} = 0 \quad \text{for } j \ne \dorder,\ 0 \le j < \dorder + \aorder.
  \label{eq:tables-moments}
\end{equation}
The first moment that does not vanish gives the leading error term. The rule's value minus $f^{(\dorder)}(x)$ is
about $\pm \errc h^{\aorder} f^{(\dorder+\aorder)}(x)$, with the error constant $\errc = \abs{\mom{\dorder+\aorder}} /
((\dorder+\aorder)!\, \den)$. The last column of each table gives the sign. The amplification $\ampl = \sum_k
\abs{c_k} / \den$ multiplies the rounding of the samples: their rounding error reaches the result as about $\ampl
\epsmach \abs{f} / h^{\dorder}$. Both constants enter the error model of \cref{ch:stencils}.

The central rules use the $\aorder + \dorder - 1$ integers around 0. For odd $\dorder$ they leave out 0, whose weight
would be zero. The forward rules use $0, 1, \ldots, \aorder + \dorder - 1$. \code{Backward<P>} is \code{Forward<P>}
mirrored. Its nodes are those of \code{Forward<P>} negated, and its numerators are those of \code{Forward<P>} times
$(-1)^{\dorder}$. It has the same $\den$, $\ampl$ and $\errc$, and the sign of its leading term changes for odd
$\aorder$. It is not tabulated.

\section{Central rules}\label{sec:tables-central}

\Cref{tab:weights-central-first,tab:weights-central-second,tab:weights-central-third,tab:weights-central-fourth}
give \code{Central<P>} for $\dorder = 1$ to 4 and $\aorder = 2$ to 10. Each pair of nodes adds two orders. $\errc$
falls quickly with $\aorder$, and $\ampl$ grows slowly. That is why the central rules of high order pay off in wide
types (\cref{tab:accuracy-matrix}). The cap at $\aorder = 10$ is a design decision
(\code{docs/deriv/PHASE2_API.md}). By the model's exponent alone, a \code{Central<12>} would cut the error of $f'$ by a
factor of about \num{\valCapGainFactorCentralTwelveFirstDouble} in double, and
\num{\valCapGainFactorCentralTwelveFirstMultiprecision} in a 50-digit type (for $f''''$,
\num{\valCapGainFactorCentralTwelveFourthDouble} and \num{\valCapGainFactorCentralTwelveFourthMultiprecision}). But
its stencil would be that much wider, and $f$ must be smooth over all of it, so the library stops there. The rules
with $\aorder = 2$ are the textbook formulas: the central difference, the three-point second difference, the
four-point rule for $f'''$ and the five-point rule for $f''''$.

\begin{table*}
\caption{\code{Central<P>} for the first derivative, $f'$.\label{tab:weights-central-first}}
\weighttablestart
\tablebody{weights-central-N1}
\weighttableend
\end{table*}

\begin{table*}
\caption{\code{Central<P>} for the second derivative, $f''$.\label{tab:weights-central-second}}
\weighttablestart
\tablebody{weights-central-N2}
\weighttableend
\end{table*}

\begin{table*}
\caption{\code{Central<P>} for the third derivative, $f'''$.\label{tab:weights-central-third}}
\weighttablestart
\tablebody{weights-central-N3}
\weighttableend
\end{table*}

\begin{table*}
\caption{\code{Central<P>} for the fourth derivative, $f^{(4)}$. The largest numerator of all the tables,
6,222,216, is the centre weight of \code{Central<10>}. Every numerator fits in 24 bits, so the integer weights of every
rule are exact in \code{float}.\label{tab:weights-central-fourth}}
\weighttablestart
\tablebody{weights-central-N4}
\weighttableend
\end{table*}

\section{Forward rules}\label{sec:tables-forward}

\Cref{tab:weights-forward-first,tab:weights-forward-second,tab:weights-forward-third,tab:weights-forward-fourth}
give \code{Forward<P>} for $\dorder = 1$ to 4 and $\aorder = 1$ to 6. A one-sided rule gains one order per node. Both
its constants are larger than those of the central rule of the same order: $\ampl$ grows to \num{\valAmplForwardSix}
for $f'$ with $\aorder = 6$, and to \num{\valAmplForwardSixFourth} for $f''''$. The cap at $\aorder = 6$ is a design
decision (\code{docs/deriv/PHASE2_API.md}). Beyond it, by the model, the amplification grows about as fast as the
higher order pays, so the library stops there.

\begin{table*}
\caption{\code{Forward<P>} for the first derivative, $f'$. For \code{Backward<P>}, negate the nodes and
multiply the numerators by $-1$.\label{tab:weights-forward-first}}
\weighttablestart
\tablebody{weights-forward-N1}
\weighttableend
\end{table*}

\begin{table*}
\caption{\code{Forward<P>} for the second derivative, $f''$. For \code{Backward<P>}, negate the nodes and
keep the numerators.\label{tab:weights-forward-second}}
\weighttablestart
\tablebody{weights-forward-N2}
\weighttableend
\end{table*}

\begin{table*}
\caption{\code{Forward<P>} for the third derivative, $f'''$. For \code{Backward<P>}, negate the nodes and
multiply the numerators by $-1$.\label{tab:weights-forward-third}}
\weighttablestart
\tablebody{weights-forward-N3}
\weighttableend
\end{table*}

\begin{table*}
\caption{\code{Forward<P>} for the fourth derivative, $f^{(4)}$. For \code{Backward<P>}, negate the nodes and
keep the numerators.\label{tab:weights-forward-fourth}}
\weighttablestart
\tablebody{weights-forward-N4}
\weighttableend
\end{table*}

\section{Default steps and model accuracy}\label{sec:tables-steps}

\Cref{tab:default-steps} gives, for every stencil and order, the exponent of the default step
$h = \epsmach^{1/(\aorder+\dorder)} \max(\abs{x}, 1)$ and the step itself for \code{double} at $\abs{x} \le 1$. For a
larger $\abs{x}$, multiply the step by $\abs{x}$. The table also gives the model's relative error
$\epsmach^{\aorder/(\aorder+\dorder)}$ at that step, and the number of evaluations. The steps are
\code{nxx::math::root_epsilon<P + N, double>()}, and \code{default_step<M, N>(x)} returns them times the scale. The
model error drops the constants $\errc$ and $\ampl$, and assumes that $f$ varies on the scale $\max(\abs{x}, 1)$,
that is, $\abs{f} \sim s^{\dorder} \abs{f^{(\dorder)}}$. It is a rule of thumb; the measured accuracies are in
\cref{tab:accuracy-matrix}.

\begin{table}
\sidecaption{Default steps, model accuracy and cost, per stencil and derivative order, for \code{double}. The step is
the one at $\abs{x} \le 1$, and the model error is relative.\label{tab:default-steps}}
\centering\footnotesize
\begin{tabular}{@{}lrrcccr@{}}
\toprule
Method & $P$ & $N$ & Exponent & Step & Model error & Evaluations\\
\midrule
\tablebody{default-steps}
\bottomrule
\end{tabular}
\end{table}

\section{Ridders' constants}\label{sec:tables-ridders}

\Cref{tab:ridders-constants} collects the constants of Ridders' method and of its one-sided forms. They are the same
for every derivative order, except where the value depends on $\dorder$. \Cref{ch:ridders} and \cref{sb:floor-coverage}
explain them, and the design record in the repository (\code{docs/deriv/PHASE3_API.md} to
\code{docs/deriv/PHASE5_API.md}) gives the evidence for each choice. The level cap, the initial-step divisors and the
one-sided evaluations per level are read from the library. The ratio, the stall and the two factors of the
untrusted-result test are literals of its code, transcribed by the generator.

The level cap counts every level attempted, restarts included. With the evaluations per level, it gives the maximum
evaluations in the last row.

\begin{table}[!htb]
\sidecaption{The constants of Ridders' method and, in the last column, of its one-sided forms
(\cref{sec:ridders-onesided}; \code{BackwardRidders} is \code{ForwardRidders} mirrored). The two differ in the column
divisor, since the error of the forward difference has every power of $h$, not only the even ones. They also differ in
the rounding floor, which the one-sided forms carry through the tableau. Where a value depends on the derivative
order, it is given for $f'$, $f''$, $f'''$ and $f''''$ in turn. The rows for the column divisor, the floor and the
restart transcribe the rules that the header documents.\label{tab:ridders-constants}}
\centering\footnotesize
\begin{tabular}{@{}>{\raggedright}p{49mm}>{\raggedright}p{28mm}>{\raggedright\arraybackslash}p{28mm}@{}}
\toprule
Constant & \code{Ridders} & \code{ForwardRidders}, \code{BackwardRidders}\\
\midrule
\tablebody{ridders-constants}
\bottomrule
\end{tabular}
\end{table}
